\documentclass[10pt,a4paper]{article}
\usepackage[margin=18mm]{geometry}
\usepackage{amsmath,amssymb,booktabs,array}
\usepackage{cite}
\renewcommand{\citepunct}{], [}
\renewcommand{\citedash}{]--[}
\input{glyphtounicode}
\pdfgentounicode=1
\usepackage[hidelinks]{hyperref}
\hypersetup{pdftitle={An urban signal-control task: model and sensitivity},
 pdfauthor={Urban traffic methods},pdfsubject={Compact methodology}}
\setlength{\parindent}{0pt}
\setlength{\parskip}{3pt}
\setlength{\abovedisplayskip}{4pt}
\setlength{\belowdisplayskip}{4pt}
\setlength{\abovedisplayshortskip}{3pt}
\setlength{\belowdisplayshortskip}{3pt}
\setlength{\emergencystretch}{2em}
\widowpenalty=10000
\clubpenalty=10000
\newcommand{\pos}[1]{[#1]_+}
\newcommand{\head}[1]{\par\vspace{4pt}\noindent\textbf{#1}\par\nobreak\vspace{1pt}}
\newenvironment{points}{\begin{list}{$\bullet$}{%
 \setlength{\leftmargin}{1.4em}\setlength{\labelwidth}{0.7em}%
 \setlength{\labelsep}{0.4em}\setlength{\itemsep}{2pt}%
 \setlength{\parsep}{0pt}\setlength{\topsep}{3pt}}}{\end{list}}
\begin{document}
\begin{center}
{\large\bfseries An urban signal-control task: model and sensitivity}
\end{center}
\vspace{-5pt}

\head{1. Research questions}
A signal can offer more left-turn service while a turner remains unable
to reach the bay. The relevant obstruction depends on which traffic is
in front, which lane it occupies, and whether another lane remains usable.
The task asks how much of that ordered information is needed for useful
signal-control sensitivities:
\begin{points}
\item \textbf{Information:} Can retained lane order explain reachable
service beyond sound aggregate models using movement counts, storage
and their own inference from traffic histories?
\item \textbf{Decision:} Does that additional information change which
effective green allocation reduces total time for the complete demand?
\item \textbf{Gradient:} Does the predicted response to a small service
reallocation agree with actual changes and help choose a better control?
\end{points}
These are separate claims: a more detailed state can improve a forward
prediction without necessarily providing a more useful control direction.

\head{2. Task and global assumptions}
\textbf{Control and population.} Consider one isolated, initially symmetric
four-leg junction with a finite set of releases and fixed routes.
One constant vector $g\geq0$, $\sum_{p=1}^4g_p=1$, allocates effective
service to NS through/right, NS left, EW through/right and EW left.
Green fractions represent average service, rather than the active phase
at a particular instant. Traffic is loaded from empty; any preparation
period is the same for every candidate control. After the control window,
a common fixed allocation clears all released demand.

\textbf{Representative approach.} There are two lanes before the bay
opening and three at the stop line. Aggregate upstream road $U$ feeds
parallel local lanes $F_1$ (left/through) and $F_2$ (through/right).
At the opening, $F_1$ branches into left-only bay $B$ and through-only
continuation $T_1$; $F_2$ continues into shared through/right lane $T_2$:
\[
 U\longrightarrow\{F_1,F_2\},\qquad
 F_1\longrightarrow\{B,T_1\},\qquad F_2\longrightarrow T_2.
\]
The bay opening lies inside the local zone. Its cross-section also
separates $F_2$ from $T_2$, so the two feeders cover the same upstream
length. In the symmetric core $B,T_1,T_2$ have the same remaining length.
Every section owns distinct road space and storage.

Left traffic enters $F_1$ and must use $B$; right traffic uses $F_2,T_2$.
Through traffic can choose $F_1$ or $F_2$ when first admitted from $U$.
Once in a feeder it stays in that lane: $F_1$ through traffic continues
to $T_1$, and $F_2$ through traffic continues to $T_2$.
The admission solver compares current feeder loads and feasible entry;
it does not move resident traffic between lanes or predict a future
signal plan. A blocked $F_1$ therefore need not stop traffic already
using $F_2$. An occupied $F_1$ head still blocks traffic behind it.
The smaller single-feeder layout remains a mechanism diagnostic;
compared models always receive the same declared geometry and permissions.

\textbf{Modeling assumptions.}
\begin{points}
\item \textbf{MA1: continuous quantity, fixed routes.} Traffic mass is
continuous vehicle-equivalent quantity. Routes, movements, lane permissions
and topology are fixed metadata. Each realized lane retains an ordered
record. Its cumulative-mass coordinate measures order, not physical metres.
\item \textbf{MA2: average signal service.} Green fractions change movement
service and downstream sharing. Road capacities, travel times, backward
wave speeds, storage and saturation parameters are fixed. The model does
not resolve red/green switching within a cycle.
\item \textbf{MA3: physical access.} An unassigned input list determines
which packet the allocator examines first. A packet already occupying a
lane also has a physical position: traffic behind a blocked head cannot
pass. Independent movement requires an independently reachable lane.
\item \textbf{MA4: entry establishes lane occupancy.} Each packet is
allocated using current receiving-lane loads and known physical limits.
Later packets see earlier actual admissions. Unadmitted mass remains
where it is, with no binding future-lane target and with priority over
later packets of its movement.
\item \textbf{MA5: causal timesteps.} Sending and receiving limits use
histories available at the start of the update. New inflow cannot leave
the same section in that update, and current departures do not instantly
create entrance room. Boundary histories are interpolated linearly.
\item \textbf{MA6: complete demand accounting.} Releases, routes and
preparation are fixed across controls. Source waiting, moving traffic and
queues all contribute until every released trip clears. Queue differences
continue to matter after the control window ends.
\end{points}

\begin{samepage}
\textbf{Travel and storage.} For each section $e$, capacity $C_e$ is a
mass rate and $K_e$ is maximum stored mass. Free-flow travel time
$\tau^f_e$ is the time to cross the section downstream. Backward-wave
time $\tau^b_e$ is the time for a congestion or space-release effect to
cross upstream. The triangular fundamental diagram relates these
quantities, and the timestep must respect both travel directions:
\begin{equation}
 K_e=C_e(\tau^f_e+\tau^b_e),\qquad
 0<\Delta t\leq\min_e\{\tau^f_e,\tau^b_e\}.
 \label{eq:compact-domain}
\end{equation}
For example, downstream departures release space, but the entrance can
use that space only after the backward wave reaches it. Numerical
lengths, demands and clearance horizons must satisfy this domain;
they are not calibrated by the equations alone.
\par\end{samepage}

\head{3. Traffic model}
The model follows a traffic journey. First determine which resident
mass has travelled far enough and where new mass can fit. Then assign
accessible traffic to lanes. Stop-line heads use their movement's service
and compete for downstream room. Finally update each physical section
and retain the histories needed for the next timestep.
Cumulative propagation and ordered route bookkeeping provide the
backbone~\cite{m0,m1}; the lane allocation and access rules below specify
the local urban interaction.

\textbf{3.1. Propagation and eligible arrivals.}
A section can send traffic only after it has travelled far enough, and
can receive traffic only when room is available at its entrance.
At $t_k=k\Delta t$, let $N^+_{e,k}$ count all earlier entries and
$N^-_{e,k}$ all earlier departures. Their difference
$n_{e,k}=N^+_{e,k}-N^-_{e,k}$ is resident mass.
The functions ${\cal H}^{\pm}_e$ read those cumulative histories between
stored boundaries by linear interpolation:
${\cal H}^{\pm}_e((k+\theta)\Delta t)
=(1-\theta)N^\pm_{e,k}+\theta N^\pm_{e,k+1}$.
Histories are zero before empty loading.

The sending calculation first counts entries early enough to reach the
exit by $t_{k+1}$ under free flow, then subtracts mass that has already
left. This gives eligible resident mass $A$. Only capacity $C_e\Delta t$
can cross during the update, so sending $S$ is the smaller amount.
Receiving uses the corresponding space account: start with storage
$K_e$, add departures whose space-release effect can reach the entrance
by $t_{k+1}$, and subtract all previous entries. The LTM limits are~\cite{m0}
\begin{align}
 A_{e,k}&=\pos{{\cal H}^+_e(t_{k+1}-\tau^f_e)-N^-_{e,k}},
 &S_{e,k}&=\min(C_e\Delta t,A_{e,k}),\label{eq:compact-sending}\\
 R_{e,k}&=\min\{C_e\Delta t,\,
 \pos{K_e+{\cal H}^-_e(t_{k+1}-\tau^b_e)-N^+_{e,k}}\}.
 \label{eq:compact-receiving}
\end{align}
Here $[z]_+=\max(z,0)$. In the first line, the delayed entrance history
sets travel-time eligibility; in the second, the delayed exit history
sets entrance-accessible room. Both $S$ and $R$ are mass allowances for
this update, not instantaneous flow rates or guaranteed actual transfers.
Signals, access and neighbouring sections can reduce realized flow.
The timestep condition makes both delayed queries no later than $t_k$.

\textbf{3.2. Lane allocation and physical access.}
The allocation problem is to admit as much of the current ready packet
as feasible, then distribute that mass toward the less loaded legal lanes.
Existing traffic is left in place. A packet is a homogeneous portion of
continuous mass carrying a route reference, current movement and original
entry/travel history~\cite{m1}.

For this packet, $q$ is its eligible mass and $s$ the sending owner's
remaining transfer allowance. For each legal, reachable receiving lane,
$u_\ell$ is its remaining admission allowance and $h_\ell$ its current
load estimate: start-of-step resident mass plus earlier actual admissions.
Thus $V$ answers ``how much can enter in total?'' It is limited by the
packet, the source and the combined receiving room. The lane solver then
chooses a common filling level $\lambda$ and additions $x_\ell$:
\begin{equation}
 V=\min(q,s,\sum_\ell u_\ell),\qquad
 x_\ell=\min\{u_\ell,[\lambda-h_\ell]_+\},\quad
 \sum_\ell x_\ell=V .
 \label{eq:compact-allocation}
\end{equation}
The unknown $\lambda$ is a queue-load level, measured in vehicle-equivalent
mass. The gap $\lambda-h_\ell$ is how much that lane would need to reach
the level. Add nothing to a lane already above it, and never exceed its
individual bound $u_\ell$. Choose the level so that all additions sum to
$V$. This is the specified ideal lane-balancing rule, not a stop-line
service cap. With initial loads $0.2,0$, admitted mass $V=0.6$ and
nonbinding lane bounds, $\lambda=0.4$ gives additions $0.2,0.4$.
A capped lane can remain below the common level. The additions are unique;
$\lambda$ itself can be nonunique when every usable lane is capped.

In the representative geometry, through packets have two alternatives
at $U\to\{F_1,F_2\}$. Left and right packets have one legal feeder each,
and movement fixes the branch at the bay opening. The same equation
therefore covers both responsive choice and forced entry. Lane-entry
room is separate from stop-line service: zero green does not immediately
prevent a vehicle from joining a lane that it can reach and that has room.

Only admitted amounts establish new lane occupancy. Update the receiving
loads by $h\leftarrow h+x$ and subtract those transfers from the remaining
sending and receiving allowances. The remainder $q-V$ stays at its current
location without a binding future-lane target. Inspect packets in recorded
order; if one movement remains partly unserved, its later packets wait,
while other movements with physical access can still be examined.

\begin{samepage}
The sending allowance limits total mass actually admitted, not how many
ready packets may be inspected. For example, an unassigned pool may list
one blocked left-turn unit before a through unit with an independently
reachable lane. With sending allowance one, inspecting only the first
unit would miss the feasible through admission. Examining the through
packet uses the same one-unit limit without moving it past a physical
blocker. Unserved mass keeps its original entry/travel history and its
relative position among the remaining entries; waiting does not restart
its travel delay.
\par\end{samepage}

Physical lane order remains binding. A residual at the front of $F_1$
prevents following traffic from reaching either $B$ or $T_1$, even if
the allocator examines its record. Traffic already in $F_2$ has its own
path to $T_2$. No traffic can jump from an occupied feeder into the
parallel lane or enter $B,T_1$ directly from $U$.

\textbf{3.3. Stop-line service and downstream sharing.}
Once a packet reaches the head of a stop-line lane, its movement sets
how quickly it can discharge. Saturation rate $s_{\ell m}$ times green
fraction $g_{p(m)}$ gives its average signal service, bounded by the
physical lane capacity $C_\ell$:
\begin{equation}
 g_p\geq0,\qquad \sum_{p=1}^4g_p=1,\qquad
 \mu_{\ell m}(g)=\min\{C_\ell,s_{\ell m}g_{p(m)}\}.
 \label{eq:compact-action}
\end{equation}
For a current ready head set $c_\ell=\mu_{\ell m}$.
Available service can still be reduced by the receiving road. Represent
output $j$'s timestep receiving allowance by the constant average rate
$\rho_j=R_{j,k}/\Delta t$. If several lanes compete for that output,
share first between inbound legs in proportion to their relevant green
fractions, then equally between ready lanes of the same leg. A lane or
leg unable to use its share releases that room to the other competitors.

Let ${\cal L}_{ij}$ contain current ready heads with positive remaining
sending, from leg $i$ to output $j$. Their combined cap is $C_{ij}$, their leg weight is $w_{ij}$,
and $W_j$ is the total rate that can be used. Common levels $\theta_j$
and $\eta_{ij}$ are chosen to satisfy the respective rate totals:
\begin{align}
 C_{ij}&=\sum_{\ell\in{\cal L}_{ij}}c_\ell,\quad
 w_{ij}=g_{p(i,j)},\quad W_j=\min(\rho_j,\sum_iC_{ij}),\nonumber\\
 r_{ij}&=\min(C_{ij},\theta_jw_{ij}),\quad
                    \sum_i r_{ij}=W_j,\label{eq:compact-leg-sharing}\\
 r_\ell&=\min(c_\ell,\eta_{ij}),\quad
                    \sum_{\ell\in{\cal L}_{ij}}r_\ell=r_{ij}.
 \label{eq:compact-lane-sharing}
\end{align}
All quantities in these sharing equations are rates. The first level
scales the green weights; the second fills uncapped lane rates equally.
This uses the same capped-filling idea as lane allocation, but allocates
\emph{discharge rate} rather than new lane mass. Zero-cap heads do not
compete, and zero receiving gives zero rates. Counting the green weight
once per leg prevents an artificial priority gain from subdividing lanes.
Shares apply to the heads ready now; later heads join when exposed.

Order determines when another movement becomes the head. Hold current
rates until a head clears, a sending allowance is exhausted, or the
update ends. With elapsed time $\sigma$, remaining head mass $q_\ell$
and remaining sending $S_\ell^{\rm rem}$, that next interval is
\begin{equation}
 \delta=\min\left\{\Delta t-\sigma,\
 \min_{\ell:r_\ell>0}
       \frac{\min(q_\ell,S_\ell^{\rm rem})}{r_\ell}\right\}.
 \label{eq:compact-head-time}
\end{equation}
During this interval, lane $\ell$ transfers $r_\ell\delta$. Then expose
new heads and recompute rates, processing simultaneous events together.
A zero-rate head blocks its own lane. This is an internal calculation of
the timestep total; it neither introduces an active signal phase clock
nor replaces the linearly interpolated propagation histories.

\textbf{3.4. Mass conservation and travel histories.}
Every internal transfer must leave one physical section and enter the
next with the same mass and route identity. If $b_{e,k}$ enters and
$y_{e,k}$ leaves section $e$, update both cumulative counts and occupancy:
\begin{equation}
 N^+_{e,k+1}=N^+_{e,k}+b_{e,k},\quad
 N^-_{e,k+1}=N^-_{e,k}+y_{e,k},\quad
 n_{e,k+1}=n_{e,k}+b_{e,k}-y_{e,k}.
 \label{eq:compact-update}
\end{equation}
The source first adds released mass $D_{o,k}$ to its waiting mass $Q_{o,k}$;
admission then moves it onto the road. Sink exit removes it from the
road population. The bounds $b\leq R\leq K-n$ and $y\leq S\leq n$
keep occupancy between zero and storage. Unadmitted mass is counted
only where it physically remains. New admissions are appended after
resident service has been selected.

\begin{samepage}
The next update also needs to know when admitted mass entered. Rather
than reconstruct a microscopic trajectory, spread the total entrance
$b>0$ uniformly across the timestep in its actual admission order.
A child packet of mass $q'$ starting at entrance rank $r$ is assigned
this entry-time interval~\cite{m1}:
\begin{equation}
 \left[t_k+\Delta t\,\frac r b,\
       t_k+\Delta t\,\frac{r+q'}b\right],
 \qquad\text{eligibility shifted by }\tau^f_e .
 \label{eq:compact-entry}
\end{equation}
Here $r$ is cumulative admitted mass before the child. The fractions
$r/b$ and $(r+q')/b$ locate its front and tail within the update; adding
free-flow time gives its earliest exit times. No entry segment exists
when inflow is zero. Internal section crossings retain route position;
the junction crossing advances it. Exact-compatible adjacent records
may merge while preserving route, order and timing. A movement grid
summarizes mass along that order without replacing the exact head record.
\par\end{samepage}

\head{4. How the task tests information and decisions}
The question is whether two states with the same aggregate quantities
can respond differently because their traffic is ordered differently.
For a local bay-opening example, take ready $F_1$ traffic
$[T(0.4),L(0.6)]$, sending allowance $1$, and receiving allowances
$R_{T_1}=0.2$, $R_B=1$. Only $0.2$ through mass enters: its residual
blocks left access. Reverse the order, keeping all lane loads and
movement totals unchanged, and $0.6$ left mass enters $B$ before $0.2$
through mass enters $T_1$. Meanwhile an independently ready $0.3$ in
$F_2$ can enter $T_2$ in either case. Total entry across the opening is
therefore $0.5$ versus $1.1$: the obstruction is selective, while the
parallel lane remains usable. This is a local arithmetic contrast,
not yet proof of different full-trip time or failure of aggregate inference.

The control objective counts everyone still travelling or waiting, not
only those who finish early. At each update charge one timestep to all
available source mass and all road mass. Summing through a common
horizon $H$ beyond complete clearance gives total system time~\cite{m0}:
\begin{equation}
 J_{\Delta t}
 =\Delta t\sum_{k=0}^{H-1}
 \left[\sum_o(Q_{o,k}+D_{o,k})+\sum_en_{e,k}\right].
 \label{eq:compact-tst}
\end{equation}
Release $D_{o,k}$ is available before that update's admission, so it is
included in that update's source charge. Transfers between road sections
do not remove anyone from the population. This is a timestep quadrature;
a nonclearing trip cannot disappear from a completed-trip-only score.

The scenario families separate simpler explanations before testing their
interaction:
\begin{center}
\setlength{\tabcolsep}{5pt}
\begin{tabular}{@{}p{0.07\linewidth}p{0.86\linewidth}@{}}
\toprule
Case & Traffic contrast\\
\midrule
V0 & Movement-exclusive paths: matched propagation and service.\\
V1 & Shared lane: head-order effects without a finite bay.\\
V2 & Finite bay: storage with order-sensitive obstruction suppressed.\\
V3 & The same bay geometry and permissions: order and access interact.\\
V4 & Held-out variation of the declared families and conditions.\\
\bottomrule
\end{tabular}
\end{center}
V2/V3 use the representative parallel-lane geometry. The smaller
single-feeder case can isolate the mechanism, but cannot establish how
an alternative lane moderates it. V0 equivalence requires matched
service and nonbinding or equivalent receiving. A different propagation
law or an unsuccessful storage model must be understood before attributing
a V3 advantage to retained order. The shared through/right lane alone
does not guarantee an order effect: its movements share green, and
become service-distinct only when their effective caps or receiving differ.

Compare B0 aggregate, B1 structured movement queues, B2 responsive
lane/storage aggregates and B3 selective ordered state under the same
population, geometry, action and time account. These roles describe
models, not a universal hierarchy of information. A macro model may
read the original fine time series through its own converter or estimator.
Its retained variables, memory, observation schedule and calibration
must be declared. Successful aggregate inference is a valid explanation;
an estimator retaining the full ordered ledger uses mesoscopic state.
Equal current counts do not imply equal information in all histories.
Distinguish a one-step test conditioned on observations, prediction with
later observations, and forecasting without them; state which data remain
fixed under a proposed service change.

The working B1 source is the link queue model~\cite{wei}; the working
B2 direction combines lane-exit allocation~\cite{gong} with a sourced
finite-storage/access adaptation, for which~\cite{liu} is a candidate.
Each complete baseline must be checked for physical soundness and task
fit. Citation alone does not establish a valid bay-entry composition.

First test the same prescribed feasible service perturbations in all
models. Then distinguish the value of the traffic representation from
the value of an optimization algorithm by matching optimization budgets.
A richer reference simulator uses one fixed mapping from effective
service to executable signals, matched releases/routes, fixed calibration
and predeclared seeds. Its population time includes insertion delay and
travel duration~\cite{sumo}, including all unfinished demand in the account.
Correct internal derivatives and a falling internal objective are
insufficient to establish useful traffic-control guidance.

\head{5. Complete sensitivity}
Start with the traffic question: if a little service is moved from one
group to another, which heads discharge more, which queues remain, and
who can reach the bay next? The start-of-step state is already fixed;
service first changes discharge and the next state. Queue and history
changes then affect later admission, spillback and total time.

Consider a feasible change $g+\epsilon d$, with $\sum_p d_p=0$.
A dot denotes its derivative at $\epsilon=0$, holding road parameters,
routes and permissions fixed. The expressions below apply while the
same heads and constraints remain active. A change in the blocking head
or binding capacity can create a kink requiring one-sided treatment.

\textbf{Lane-allocation response.} A previous service change can alter
both lane loads and the room available for entry. Among lanes receiving
an uncapped positive addition, the common filling level must adjust so
the additions still total $V$. Let ${\cal A}$ be these free lanes and
${\cal U}$ the lanes at their upper admission bounds. Differentiating
the mass balance gives
\begin{align}
 \lambda&=\frac{V-\sum_{\cal U}u_\ell+\sum_{\cal A}h_\ell}
                  {|{\cal A}|},\qquad
 \dot\lambda=\frac{\dot V-\sum_{\cal U}\dot u_\ell+
                                  \sum_{\cal A}\dot h_\ell}
                  {|{\cal A}|},\nonumber\\
 \dot x_\ell&=
 \begin{cases}
 \dot\lambda-\dot h_\ell,&\ell\in{\cal A},\\
 \dot u_\ell,&\ell\in{\cal U},\\
 0,&\text{inactive lower branch}.
 \end{cases}
 \label{eq:compact-allocation-derivative}
\end{align}
An uncapped lane responds through the common level minus its own load
change; a capped lane follows its available room. Differentiate $V$
through its limiting argument, then carry the additions into receiving
loads, remaining allowances and retained mass
$\dot q^{\rm rem}=\dot q-\dot V$ before processing the next packet.
With no free lanes, use the cap or stable zero branches directly.
A zero current admission does not erase the response of waiting mass.
At a fixed pre-state this admission law has no direct green input;
its control response arrives through the preceding traffic evolution.

\textbf{Service and shared receiving.} Green directly changes a head's
service cap: $\dot\mu_{\ell m}=s_{\ell m}d_{p(m)}$ when signal-limited,
and zero when physical capacity limits it. It also changes the leg's
sharing weight. Receiving changes through previous road states, with
$\dot\rho_j=\dot R_{j,k}/\Delta t$. The shared rate balance must reflect
all of these changes. For uncapped legs ${\cal F}$ and capped legs
${\cal C}$, its response is
\begin{equation}
 \dot\theta=
 \frac{\dot W-\sum_{\cal C}\dot C_{ij}
               -\theta\sum_{\cal F}\dot w_{ij}}
      {\sum_{\cal F}w_{ij}},\qquad
 \dot r_{ij}=
 \begin{cases}
 \dot C_{ij},&i\in{\cal C},\\
 w_{ij}\dot\theta+\theta\dot w_{ij},&i\in{\cal F}.
 \end{cases}
 \label{eq:compact-sharing-derivative}
\end{equation}
Within a leg use the corresponding equal-share response. If no competitor
is uncapped, follow cap branches without dividing by an empty weight sum.

Faster service can also expose the next head sooner. For an interior
interval $\delta=q/r$, its time response is
$\dot\delta=\dot q/r-q\dot r/r^2$. Total discharge changes through both
the rates and the durations for which they apply:
\begin{equation}
 y=\sum_a r_a\delta_a,\qquad
 \dot y=\sum_a(\dot r_a\delta_a+r_a\dot\delta_a).
 \label{eq:compact-service-derivative}
\end{equation}
Index $a$ runs over internal service intervals. The final interval changes
by the negative sum of the earlier duration changes so the whole update
still lasts $\Delta t$. Freezing head-clearance times would omit this
physical route by which one movement's service affects another.

For illustration, V1's shared left/through diagnostic has a ready left
block $q_L$ before sufficient through traffic. With nonbinding receiving
and physical caps, left service first clears that block; through service
uses only the remaining time:
\begin{equation}
 y_T=s_Tg_T\left(\Delta t-\frac{q_L}{s_Lg_L}\right),
 \qquad 0<\frac{q_L}{s_Lg_L}<\Delta t.
 \label{eq:compact-service}
\end{equation}
With unit timestep and saturations, $q_L=0.2$ and
$(g_L,g_T)=(0.4,0.6)$ give through discharge $0.3$. Moving a small fraction
from through to left has through-discharge derivative $0.25$ on this
branch: earlier head clearance outweighs the reduced through rate.
This is a V1 service illustration, not left-turn permission in the
representative through-only $T_1$, nor a claim about complete-time benefit.

\textbf{Arrival timing and the full state.} Earlier discharge also changes
boundary counts and the entry times assigned to subsequent traffic.
Differentiate the delayed histories and the reconstruction, so that
later readiness changes with both packet mass and arrival time.
For a positive-inflow branch the front-time response is
\begin{equation}
 \dot t_{\rm front}
 =\Delta t\left(\frac{\dot r}{b}-\frac{r\dot b}{b^2}\right),
 \label{eq:compact-entry-derivative}
\end{equation}
The tail has the analogous response. If fraction $f$ of a packet's
entry/travel segment is ready, ready mass responds as
$f\dot q+q\dot f$. Zero inflow creates no timing segment.

Collect road masses, histories, ordered records and sources in state
$X_k$, and write $P_k=\dot X_k$. If $u_k$ is the green allocation applied
in update $k$, the chain rule propagates the effect of a service change:
\begin{equation}
 X_{k+1}=F_k(X_k,u_k),\qquad
 P_{k+1}=F_{k,X}P_k+F_{k,u}\dot u_k .
 \label{eq:compact-chain}
\end{equation}
The first term carries earlier state differences forward; the second
adds the current control change. Preparation is common and has zero
candidate response. During control $\dot u_k=d$; after switching to fixed
clearance $\dot u_k=0$, but $P_k$ continues to evolve. Switching the knob
back does not undo the queue differences it created.

Finally, all released mass is either still waiting/travelling or has
exited. With cumulative releases $D_{\leq k}$ through update $k$ and
sink exits $Z_k$ before that update, the same objective and its response
can be written
\begin{equation}
 J_{\Delta t}=\Delta t\sum_{k=0}^{H-1}(D_{\leq k}-Z_k),
 \qquad
 \dot J_{\Delta t}=-\Delta t\sum_{k=0}^{H-1}\dot Z_k .
 \label{eq:compact-objective-derivative}
\end{equation}
The sign is causal: earlier additional exits reduce population time.
Use a common horizon that clears the compared control neighbourhood.
This complete sensitivity is distinct from proving its numerical accuracy
or downstream usefulness. Implementation, numerical derivative checks,
baseline qualification and richer-reference tests remain necessary evidence.
The \href{urban_methodology_specification.pdf}{mathematical companion}
gives full boundary cases and the derivation; source/adaptation records
identify the origin of each operator.

\makeatletter
\renewcommand{\section}{\@ifstar{\head}{\head}}
\makeatother
\begin{thebibliography}{6}
\setlength{\itemsep}{0pt}\setlength{\parskip}{0pt}
\bibitem{m0} Cumulative-count LTM and population-time discretization,
repository methods specification:
\href{../../../../docs/milestones/milestone_0_ltm_baseline.md}{milestone plan}.
Its recorded original basis is Yperman (2007) and UNsim v0.12.0;
the equations here use the verified repository transcription.
\bibitem{m1} Ordered route-mass representation, conservation and affine timing,
repository methods specification:
\href{../../../../docs/milestones/milestone_1_ordered_route_mass_ledger.md}
{milestone plan}. Urban pool selection/access is the explicit adaptation
defined in this report.
\bibitem{wei} Wei et al., link queue model,
\href{https://arxiv.org/html/2310.12249v2}{arXiv:2310.12249v2};
journal version doi:
\href{https://doi.org/10.1016/j.trc.2024.104937}{10.1016/j.trc.2024.104937}.
\bibitem{gong} Gong et al., lane-based macroscopic node model (2025),
doi: \href{https://doi.org/10.1016/j.trb.2025.103208}
{10.1016/j.trb.2025.103208}.
\bibitem{liu} Liu and Chang, lane-group storage/blockage model,
\emph{Transportation Research C}, vol. 19, pp. 130--144, 2011,
\href{https://attap.umd.edu/wp-content/uploads/2022/05/TRC522.pdf}
{author-hosted article}; doi: 10.1016/j.trc.2010.04.005.
\bibitem{sumo} SUMO official documentation,
\href{https://sumo.dlr.de/docs/Simulation/Traffic_Lights.html}
{traffic-light programs} and
\href{https://sumo.dlr.de/docs/Simulation/Output/TripInfo.html}
{TripInfo population-time fields}.
\end{thebibliography}
\end{document}
